\section{Analysis and Results}

Our evaluation addresses three questions:
\textbf{(1)} Can video-derived supervision enable reliable goal reaching and obstacle avoidance on a physical robot?
\textbf{(2)} How does VEGA compare with existing VLAs on trajectory quality and multimodal goal reaching?
\textbf{(3)} How reliable is the generated geometric supervision, and how robust are VEGA's performance gains to the geometry source and collision criteria?

\begin{figure*}[t]
    \centering

    % ==================== Left: Real-world qualitative results ====================
    \begin{minipage}[t]{0.57\textwidth}
        \vspace{0pt}
        \centering

        \includegraphics[width=\linewidth]{fig/vega_comps.jpg}

        \captionof{figure}{
        \textbf{Real-world goal-conditioned navigation with VEGA.}
        Qualitative comparison across three real-world scenarios with static and dynamic obstacles.
        Each trajectory is directed toward the yellow-star goal; check marks indicate success and X marks indicate collisions or failures.
        VEGA reaches the goal while avoiding obstacles, whereas baseline VLAs often collide, stop short, or miss the target.
        }
        \label{fig:comps}
    \end{minipage}
    \hfill
    % ==================== Middle: Isaac Sim scenarios ====================
    \begin{minipage}[t]{0.19\textwidth}
        \vspace{0pt}
        \centering

        \includegraphics[width=\linewidth]{fig/sim_fig.jpg}

        \captionof{figure}{
        \textbf{Isaac Sim scenarios.}
        Robot-view snapshots for Scenarios~1--3; the stationary person in white is the goal.
        }
        \label{fig:isaac_scenarios}
    \end{minipage}
    \hfill
    % ==================== Right: Isaac Sim quantitative results ====================
    \begin{minipage}[t]{0.22\textwidth}
        \vspace{0pt}
        \centering

        \scriptsize
        \setlength{\tabcolsep}{1.5pt}
        \renewcommand{\arraystretch}{1.22}

        \begin{tabular}{llrrr}
        \toprule
        Scene & Method & S$\uparrow$ & C$\downarrow$ & Cl$\uparrow$ \\
        \midrule
        Scn1 & VEGA & \textbf{1.00} & \textbf{0.00} & \textbf{0.85} \\
           & OmniVLA     & 0.42 & 0.58 & 0.04 \\
           & NaVILA      & 0.80 & \textbf{0.00} & 0.84 \\
           & TIC-VLA     & 0.18 & 0.50 & 0.45 \\
        \midrule
        Scn2 & VEGA & \textbf{0.90} & \textbf{0.10} & 0.49 \\
           & OmniVLA     & 0.06 & 0.96 & 0.00 \\
           & NaVILA      & \textbf{0.90} & 0.02 & \textbf{0.53} \\
           & TIC-VLA     & \textbf{0.90} & 0.04 & 0.48 \\
        \midrule
        Scn3 & VEGA & \textbf{1.00} & \textbf{0.00} & 1.06 \\
           & OmniVLA     & 0.04 & 0.16 & 0.13 \\
           & NaVILA      & 0.00 & 0.00 & 1.52 \\
           & TIC-VLA     & 0.74 & 0.06 & 0.61 \\
        \bottomrule
        \end{tabular}

        \captionof{table}{
        Isaac Sim results over 50 trials per method and scenario.
        $S$ and $C$ are rates; $Cl$ is clearance in meters. \ours{} outperforms baselines overall in 50 simulations per scenario.
        }
        \label{tab:isaac_sim}
    \end{minipage}

    \vspace{-4ex}
\end{figure*}
\subsection{Experimental Setup}

\noindent \textbf{Baselines.}
We compare VEGA with OmniVLA~\cite{hirose2025omnivlaomnimodalvisionlanguageactionmodel}, TIC-VLA~\cite{huang2026tic}, and NaVILA~\cite{cheng2025navilaleggedrobotvisionlanguageaction}, three state-of-the-art navigation VLAs.
This comparison assesses whether geometric supervision derived from unlabeled Internet videos enables effective goal-directed navigation and obstacle avoidance relative to VLAs trained using robot demonstrations.
Goal specifications are provided through each model's supported interface while preserving the intended target across methods.
For real-world execution, policy outputs are passed to a shared control stack~\cite{seneviratne2026chop} that provides no additional path planning or obstacle avoidance, but allows all methods to replan from updated observations before completing a prediction.
NaVILA's high-level navigation commands are converted into equivalent local waypoints and executed through this shared stack.
We do not use NaVILA's native LiDAR-conditioned obstacle-avoiding locomotion policy, since our goal is to evaluate obstacle avoidance encoded in the VLA predictions themselves rather than corrections introduced by downstream navigation modules.
Thus, all methods are evaluated under a common execution interface without additional obstacle avoidance.

\paragraph{Implementation Details.}
During dataset construction, we use a $25\,\mathrm{m}\times20\,\mathrm{m}$ BEV at $0.02\,\mathrm{m}$/cell resolution, centered on the robot's forward workspace. We use a safety distance of $d_{\mathrm{safe}}=0.5\,\mathrm{m}$, corresponding to the robot radius, and generate waypoint chunks with horizon $H=50$. Goal reaching is defined as terminating within 0.5\,m of the target.

\noindent \textbf{Metrics.}
We evaluate each method along three axes:
\begin{itemize}[noitemsep,topsep=1pt]
    \item \textbf{Goal reaching.}
    On VEGA-Bench, we report average normalized goal progress, defined as the fractional reduction in goal distance normalized by the initial robot-goal distance.
    In real-world trials, success requires the robot to reach within a threshold distance of the intended goal without any collision or safety intervention.

    \item \textbf{Collision avoidance.}
    On VEGA-Bench, collision rate is the fraction of trajectories that enter occupied, unknown, or out-of-bounds ESDF regions.
    In real-world trials, collision rate is the fraction of trials with physical contact or safety intervention.

    \item \textbf{Obstacle clearance.}
    We report the average minimum distance to the nearest obstacle along non-colliding trajectories.
\end{itemize}

\begin{table}[t]
\vspace{2ex}
\centering
% \footnotesize
\setlength{\tabcolsep}{3pt}
\renewcommand{\arraystretch}{1.1}

\resizebox{\columnwidth}{!}{%
\begin{tabular}{@{}l*{12}{c}@{}}
\toprule
& \multicolumn{3}{c}{Scenario 1}
& \multicolumn{3}{c}{Scenario 2}
& \multicolumn{3}{c}{Scenario 3}
& \multicolumn{3}{c}{Scenario 4} \\
\cmidrule(lr){2-4}
\cmidrule(lr){5-7}
\cmidrule(lr){8-10}
\cmidrule(lr){11-13}
Method
& S$\uparrow$ & C$\downarrow$ & Cl$\uparrow$
& S$\uparrow$ & C$\downarrow$ & Cl$\uparrow$
& S$\uparrow$ & C$\downarrow$ & Cl$\uparrow$
& S$\uparrow$ & C$\downarrow$ & Cl$\uparrow$ \\
\midrule
OmniVLA
& 0.15 & 0.85 & 0.35
& 0.0 & 0.65 & 0.32
& 0.0 & 0.5 & 0.44
& 0 & 1.0 & N/A \\
TIC-VLA
& 0.0 & 1.0 & N/A
& 0.1 & 0.85 & 0.37
& 0.0 & 1.0 & N/A
& 0 & 1.0 & N/A \\
NaVILA
& 0.0 & 1.0 & N/A
& 0.0 & 0.8 & 0.35
& 0.0 & 0.45 & 0.31
& 0 & 0.75 & 0.25 \\
\textbf{VEGA}
& \textbf{0.95} & \textbf{0.05} & \textbf{0.93}
& \textbf{0.75} & \textbf{0.25} & \textbf{0.56}
& \textbf{1.0} & \textbf{0.0} & \textbf{0.94}
& \textbf{0.95} & \textbf{0.05} & \textbf{0.52} \\
\bottomrule
\end{tabular}%
}

\caption{Real-world navigation over 20 trials.
S: success rate; Cl: clearance in meters, with N/A when
all trials collide; C: average collision rate
for Scenarios 1--4.}
\label{tab:real_results}
\label{tab:scenario4}
\vspace{-2ex}
\end{table}

\begin{table}[t]
\centering
\footnotesize
\setlength{\tabcolsep}{3pt}
\renewcommand{\arraystretch}{1.1}

\resizebox{\columnwidth}{!}{%
\begin{tabular}{@{}l*{9}{c}@{}}
\toprule
& \multicolumn{3}{c}{VEGA-Bench}
& \multicolumn{3}{c}{Collision composition}
& \multicolumn{3}{c}{Goal conditioning} \\
\cmidrule(lr){2-4}
\cmidrule(lr){5-7}
\cmidrule(lr){8-10}
Method
& G$\uparrow$ & C$\downarrow$ & Cl$\uparrow$
& $C_{\mathrm{occ}}$ & $C_{\mathrm{unk}}$ & $C_{\mathrm{oob}}$
& Text$\uparrow$ & Image$\uparrow$ & WP$\uparrow$ \\
\midrule
OmniVLA
& 0.88 & 6.75 & 2.35
& 71.4\% & 23.8\% & 4.74\%
& 0.8 & 0.2 & 0.6 \\
TIC-VLA
& 0.84 & 9.01 & 1.34
& 68.1\% & 26.6\% & 5.27\%
& 0.8 & N/A & N/A \\
NaVILA
& 0.77 & 7.40 & 2.18
& 68.4\% & 26.8\% & 4.9\%
& 0.6 & N/A & N/A \\
\textbf{VEGA}
& \textbf{0.89} & \textbf{5.52} & \textbf{2.87}
& 75.9\% & 20.93\% & 3.16\%
& \textbf{1.0} & \textbf{1.0} & \textbf{0.8} \\
\bottomrule
\end{tabular}%
}

\caption{VEGA-Bench evaluation, collision composition,
and goal-reaching performance by modality.
G: goal progress; C: collision percentage;
Cl: clearance in meters.
Collision composition reports percentages of collisions
in occupied ($C_{\mathrm{occ}}$), unknown ($C_{\mathrm{unk}}$),
and out-of-bounds ($C_{\mathrm{oob}}$) cells. Goal-modality results report success over 20 real-world trials per supported modality using the Scenario~4 layout with dynamic obstacles removed.;
WP denotes waypoint goals.}
\label{tab:vega_bench_results}
\label{tab:collision_decomposition}
\label{tab:goal_modality_results}
\vspace{-5ex}
\end{table}



\subsection{Real-World Robot Evaluation}

% For real-world trials, all methods are executed through the shared ROS-based stack from CHOP~\cite{seneviratne2026chop}.
% NaVILA's native action outputs are mapped to equivalent local waypoints for execution, while the other methods retain their native navigation outputs.
% The shared stack performs no additional path planning or obstacle avoidance and allows all methods to replan from updated observations before completing a prediction.

To assess real-world transfer, we deploy the methods on a quadruped in cluttered environments with both static and dynamic obstacles, as shown in Fig.~\ref{fig:comps}.
Across four scenarios with 20 trials per scenario, VEGA achieves the highest success rate and the fewest collisions among the evaluated baselines (Table~\ref{tab:real_results}). 
These results demonstrate that video-derived trajectory supervision enables effective goal reaching and obstacle avoidance on a physical robot in the tested environments.

To separately evaluate goal conditioning from dynamic-obstacle avoidance, we conduct 20 trials per supported goal modality using the Scenario~4 layout with the dynamic obstacles removed (Table~\ref{tab:vega_bench_results}).
These trials test whether each policy can interpret the specified goal and generate goal-directed motion under a simpler navigation setting.
Across text, image, and waypoint goals, VEGA produces distinct trajectories that reach the corresponding targets, while some baselines often generate similar paths despite changes in the goal specification.

\subsection{Isaac Sim Evaluation}

We additionally evaluate navigation in three Isaac Sim scenarios containing static obstacles and moving pedestrians.
%illustrated in Fig.~\ref{fig:isaac_scenarios}. 
Trials vary pedestrian motion onset using identical timing conditions across methods.
Following the real-world evaluation, we report success rate, collision rate, and average minimum clearance over non-collision trials in Table~\ref{tab:isaac_sim}.

% % Requires \usepackage{booktabs}
% \begin{table}[t]
% \raggedleft
% \footnotesize
% \setlength{\tabcolsep}{2pt}
% \begin{tabular}{llrrr}
% \toprule
% Scene & Method & S & C & Cl \\
% \midrule
% S1 & VEGA (ours) & 100 &   0 & 0.805 \\
%    & OmniVLA     &  25 &  75 & 0.106 \\
%    & NaVILA      & 100 &   0 & 1.047 \\
%    & TIC-VLA     &  50 & 100 & --    \\
% \midrule
% S2 & VEGA (ours) & 100 &   0 & 0.599 \\
%    & OmniVLA     &   0 & 100 & --    \\
%    & NaVILA      & 100 &   0 & 0.455 \\
%    & TIC-VLA     & 100 &   0 & 0.368 \\
% \midrule
% S3 & VEGA (ours) &  75 &  25 & 0.064 \\
%    & OmniVLA     & 100 &   0 & 0.256 \\
%    & NaVILA      &   0 &  25 & 0.311 \\
%    & TIC-VLA     &  75 &  50 & -0.018 \\
% \bottomrule
% \end{tabular}
% \caption{Isaac Sim results with fifty trials per method
% per scenario. S and C are percentages. Cl denotes average
% minimum clearance over non-collision trials, in meters;
% ``--'' indicates no eligible trials.}
% \label{tab:isaac_sim}
% \vspace{-5ex}
% \end{table}
% \begin{figure}[t]
% \centering

% ==================== Left: Isaac Sim scenarios ====================
% \begin{minipage}[t]{0.5\columnwidth}
%     \vspace{0pt}
%     \centering

%     \includegraphics[width=\linewidth]{fig/sim_fig.jpg}

%     \captionof{figure}{
%     \textbf{Isaac Sim scenarios.}
%     Robot-view snapshots for Scenarios~1--3; the stationary person in white is the goal.
%     }
%     \label{fig:isaac_scenarios}
% \end{minipage}
% \hfill
% % ==================== Right: navigation results ====================
% \begin{minipage}[t]{0.45\columnwidth}
%     \vspace{0pt}
%     \centering

%     \scriptsize
%     \setlength{\tabcolsep}{1.5pt}
%     \renewcommand{\arraystretch}{1.05}

%     \begin{tabular}{llrrr}
%     \toprule
%     Scene & Method & S & C & Cl \\
%     \midrule
%     S1 & VEGA (ours) & 100 &   0 & 0.805 \\
%        & OmniVLA     &  25 &  75 & 0.106 \\
%        & NaVILA      & 100 &   0 & 1.047 \\
%        & TIC-VLA     &  50 & 100 & --    \\
%     \midrule
%     S2 & VEGA (ours) & 100 &   0 & 0.599 \\
%        & OmniVLA     &   0 & 100 & --    \\
%        & NaVILA      & 100 &   0 & 0.455 \\
%        & TIC-VLA     & 100 &   0 & 0.368 \\
%     \midrule
%     S3 & VEGA (ours) &  75 &  25 & 0.064 \\
%        & OmniVLA     & 100 &   0 & 0.256 \\
%        & NaVILA      &   0 &  25 & 0.311 \\
%        & TIC-VLA     &  75 &  50 & -0.018 \\
%     \bottomrule
%     \end{tabular}

%     \captionof{table}{
%     Isaac Sim results over 50 trials per method and scenario.
%     $S$ and $C$ are percentages; $Cl$ is clearance in meters.
%     }
%     \label{tab:isaac_sim}
% \end{minipage}

% \vspace{-3ex}
% \end{figure}

\begin{table*}[!t]
\centering
\footnotesize
\setlength{\tabcolsep}{3pt}
\renewcommand{\arraystretch}{1.1}
\resizebox{\textwidth}{!}{%
\begin{tabular}{@{}l*{12}{c}@{}}
\toprule
& \multicolumn{3}{c}{OmniVLA}
& \multicolumn{3}{c}{TIC-VLA}
& \multicolumn{3}{c}{NaVILA}
& \multicolumn{3}{c}{VEGA} \\
\cmidrule(lr){2-4}
\cmidrule(lr){5-7}
\cmidrule(lr){8-10}
\cmidrule(lr){11-13}
Pointcloud source
& $G$$\uparrow$ & $C$$\downarrow$ & $Cl$$\uparrow$
& $G$$\uparrow$ & $C$$\downarrow$ & $Cl$$\uparrow$
& $G$$\uparrow$ & $C$$\downarrow$ & $Cl$$\uparrow$
& $G$$\uparrow$ & $C$$\downarrow$ & $Cl$$\uparrow$ \\
\midrule
MoGe-2
& 0.88/1.00 & 6.75 & 2.35/1.00
& 0.84/1.00 & 9.01 & 1.34/1.00
& 0.77/1.00 & 7.40 & 2.18/1.00
& 0.89/1.00 & 5.52 & 2.87/1.00 \\
UniDepthV2
& 0.87/0.97 & 6.91 & 2.31/0.95
& 0.84/0.95 & 9.43 & 1.91/0.94
& 0.76/0.96 & 7.56 & 2.14/0.94
& 0.86/0.97 & 5.89 & 2.61/0.91 \\
UniK3D
& 0.86/0.95 & 6.58 & 2.39/0.93
& 0.80/0.94 & 8.96 & 1.97/0.92
& 0.75/0.94 & 7.23 & 2.22/0.92
& 0.85/0.95 & 5.39 & 2.82/0.95 \\
Metric3Dv2
& 0.84/0.92 & 7.09 & 2.27/0.90
& 0.76/0.90 & 9.65 & 1.82/0.89
& 0.73/0.92 & 7.79 & 2.10/0.90
& 0.83/0.92 & 5.88 & 2.64/0.92 \\
% \midrule
% SCAND (LiDAR)
% & 0.37 & 5.42 & 2.64
% & 0.31 & 7.68 & 2.19
% & 0.29 & 6.13 & 2.47
% & 0.34 & 3.61 & 3.06 \\
\bottomrule
\end{tabular}%
}
\caption{Benchmark sensitivity to pointcloud source.
For goal progress ($G$) and clearance ($Cl$), each monocular-model entry reports mean $\mu$/Spearman $\rho$ relative to MoGe-2.
Collision ($C$) reports mean collision percentage.
$G$: goal progress ($\uparrow$);
$C$: collision percentage ($\downarrow$);
$Cl$: clearance in meters ($\uparrow$).}
\label{tab:geometry_sensitivity}
\vspace{-4ex}
\end{table*}

\subsection{VEGA-Bench Evaluation}

To evaluate goal-conditioned navigation at scale, we benchmark all methods on VEGA-Bench.
For the main comparison, all methods are evaluated on the common text-conditioned subset, while modality-specific real-world evaluations are reported separately.
The VEGA-Bench test set contains approximately 12.5 hours of held-out videos, with the split performed at the city-level so that frames and their derived goals and trajectories from a given video do not appear in both training and evaluation.
Predicted trajectories are evaluated against the reconstructed scene geometry using goal progress, collision rate, and obstacle clearance.
As shown in Table~\ref{tab:vega_bench_results}, VEGA achieves competitive goal progress while reducing collision rate and increasing obstacle clearance relative to the evaluated baselines.